% Pauli/Weyl algebra targets. Detailed matrix multiplication is in the owner proof zone.
后面描述自旋 $1/2$ 场时会遇到一种具有两个复分量、并按 Lorentz 群二维复表示变换的对象；这种二分量对象称为 Weyl 旋量。Weyl 旋量所属的分量空间就是二维复线性空间。任意 $2\times2$ 复矩阵都可由单位矩阵与三个无迹厄米矩阵张成；要求这三个生成元闭合成 $SU(2)$ 转动代数后，自然选取 Pauli 矩阵
\begin{equation}
 \sigma_1=\begin{pmatrix}0&1\\1&0\end{pmatrix},\qquad
 \sigma_2=\begin{pmatrix}0&-\ii\\ \ii&0\end{pmatrix},\qquad
 \sigma_3=\begin{pmatrix}1&0\\0&-1\end{pmatrix}.
 \label{eq:pauli-explicit-r10}
\end{equation}
它们的乘法闭合为
\begin{equation}
 \boxed{
 \sigma_i\sigma_j
 =\delta_{ij}\id_2+\ii\epsilon_{ijk}\sigma_k,}
 \label{eq:pauli-product-proof-r10}
\end{equation}
因而对称与反对称部分分别满足
\begin{equation}
 \boxed{
 \{\sigma_i,\sigma_j\}=2\delta_{ij}\id_2,
 \qquad
 [\sigma_i,\sigma_j]=2\ii\epsilon_{ijk}\sigma_k.}
 \label{eq:pauli-comm-anticomm-r10}
\end{equation}

对单位向量 $\hat{\mathbf n}$，\eqnrefs{eq:pauli-product-proof-r10} 立即给出
$(\hat{\mathbf n}\cdot\boldsymbol\sigma)^2=\id_2$。因此有限转动的矩阵指数能够按奇偶幂精确求和：
\begin{equation}
 \boxed{
 \ee^{-\ii\theta\hat{\mathbf n}\cdot\boldsymbol\sigma/2}
 =\id_2\cos\frac{\theta}{2}
 -\ii\hat{\mathbf n}\cdot\boldsymbol\sigma\sin\frac{\theta}{2}.}
 \label{eq:pauli-exponential-r68}
\end{equation}
把纯虚系数 $-\ii\theta$ 换成实快度 $\pm\zeta$，同一代数同时给出 Weyl 推动中的双曲函数形式。

三维反对称张量的缩并恒等式取为
\begin{equation}
 \boxed{
 \epsilon_{ijk}\epsilon_{ilm}
 =\delta_{jl}\delta_{km}-\delta_{jm}\delta_{kl},
 \qquad
 \epsilon_{ijk}\epsilon_{ijm}=2\delta_{km}.}
 \label{eq:epsilon-contractions-r10}
\end{equation}
这些关系随后把三维自旋代数嵌入四维 Lorentz 结构。定义
\begin{equation}
 \sigma^\mu=(\id_2,\boldsymbol\sigma),
 \qquad
 \bar\sigma^\mu=(\id_2,-\boldsymbol\sigma),
 \label{eq:sigma-barsigma-foundation-r10}
\end{equation}
则 Pauli 乘法定律等价于四维交织恒等式
\begin{equation}
 \boxed{
 \sigma^\mu\bar\sigma^\nu+\sigma^\nu\bar\sigma^\mu
 =2\eta^{\mu\nu}\id_2,}
 \label{eq:sigmabar-proof-r10}
\end{equation}
交换 $\sigma$ 与 $\bar\sigma$ 的次序也成立。\eqnrefs{eq:sigmabar-proof-r10} 说明这组矩阵的对称反对易部分恰好重现 Minkowski 度规。这里出现这种关系并不是为了再引入一个新的对称群，而是为了完成一个具体的代数任务：相对论自由粒子的质量壳条件 $p^2=m^2c^2$ 对四动量是二次的；若希望以后把运动方程写成对四动量的一阶形式，就需要一组常矩阵 $\Gamma^\mu$，使一阶算符 $\Gamma^\mu p_\mu$ 的平方只留下 Lorentz 标量 $p^2$。这要求矩阵的对称乘积满足
\begin{equation*}
 \{\Gamma^\mu,\Gamma^\nu\}=2\eta^{\mu\nu}I.
\end{equation*}
由满足这组反对易关系的生成元及其有限乘积构成的结合代数，称为 Minkowski 度规对应的 \emph{Clifford 代数}。因此 Clifford 代数是一种把二次 Lorentz 不变量“线性化”的结合矩阵代数；时空对称变换仍由 Lorentz 群及其自旋表示描述。\eqnrefs{eq:sigmabar-proof-r10} 已经在两个二分量空间之间实现了同一种代数结构，后面把左右两个二分量空间合并后，就会得到作用在四分量空间上的 $\gamma^\mu$。
